Until now, spinors have been discussed as wave functions for the intrinsic
angular momentum of elementary particles.  Formally, however, the spin of a
single particle is no different from the total angular momentum of any system
regarded as a whole, when its internal structure is set aside.

\SourcePage{263}{266}
Consequently, the transformation properties of spinors apply equally to the
behaviour under spatial rotations of the wave functions $\psi_{jm}$ for any
particle (or particle system) having total angular momentum $j$, irrespective
of whether that angular momentum is orbital or intrinsic.  There must
therefore be a definite correspondence between the transformation laws of the
eigenfunctions $\psi_{jm}$ under rotations of the coordinate system and those
of the components of a symmetric spinor of rank $2j$.
